\section{Related work}\label{sec:related}

\paragraph{Thermodynamics benchmarks and agents.}
ThermoQA is a benchmark of thermodynamic questions for language models; it is a single-author
preprint \citep{duzkar2026thermoqa}. Song et al.\ built ReAct agents for gas-path analysis of
gas turbines and found that smaller models struggled with tool use and parameter extraction
\citep{song2024gaspath}. Both evaluate models; neither trains one against a simulator.

\paragraph{Design and flowsheet agents.}
Text2Sim uses multiple prompted agents, with a frozen model, to turn a text description into a
chemical-process simulation \citep{tian2026text2sim}. Li et al.\ treat thermodynamic cycle co-design
as a graph-based hierarchical reinforcement-learning problem with a learned surrogate and no language
model \citep{li2026thermocycle}. Our Design task uses a language model to choose a configuration and
continuous parameters, but the final check is the physical simulator rather than a surrogate.

\paragraph{Reinforcement learning with verifiable and solver-grounded rewards.}
GRPO \citep{shao2024deepseekmath} and its use for reasoning \citep{guo2025deepseekr1} replace a learned
value function with group-normalized rewards. Tulu~3 introduced reinforcement learning with verifiable
rewards for post-training and motivates a supervised stage first: sparse verifiable rewards need a policy
that already succeeds sometimes \citep{lambert2024tulu3}. Solver-informed reinforcement learning uses
a solver in the loop for optimization modeling \citep{chen2025sirl}. Reinforcement learning from execution feedback
in code uses test outcomes as the reward \citep{gehring2024rlef}, and the loss we use follows the DAPO
variant of GRPO \citep{yu2025dapo}. The closest methodological precedent is Luo et al., who run
supervised fine-tuning followed by GRPO with a solver-grounded reward for modeling nuclear
thermal-hydraulic systems, with LoRA rank 8 and a group of 8; they report solver convergence rising
from 72.4\% after supervised fine-tuning to 88.8\% after GRPO \citep{luo2026nuclear}. Our work follows the same
two stages with a different reward: the physical second-law checks of \cref{sec:verifiers}.

\paragraph{Parameter-efficient reinforcement learning.}
LoRA \citep{hu2021lora} is the adapter method we use. Reinforcement learning has been reported to need little adapter capacity, with a learning
rate about ten times that of full fine-tuning \citep{schulman2025lora}; our GRPO stage uses a smaller
rank than the SFT stage for that reason. Adapting a decoder-only model into an encoder-decoder model
\citep{zhang2025t5gemma} is an optional ablation only, because reinforcement-learning tooling
targets decoder-only models.

\paragraph{Data-center on-site generation, waste heat and absorption cooling.}
Combined heat and power for data centers was assessed in an Oak Ridge report \citep{darrow2009chp}.
Published measurements of the Frontier supercomputer give the load and the coolant temperatures that
set our operating points \citep{sun2024frontier,jadhav2026frontierml}. The coefficient of performance
of a single-effect absorption chiller was measured at 0.55 to 0.70 with 88\,$^\circ$C hot water,
12.5\,$^\circ$C chilled water and 31\,$^\circ$C cooling water \citep{ochoa2019absorption}; our
absorption model is a system-level fraction of the reversible COP, capped at 0.8, not a property model
(\cref{sec:limits}). Exergy analysis and its use in design follows standard texts \citep{bejan1996thermal}.
Gas properties come from the GRI-Mech 3.0 polynomials \citep{smith1999grimech} and liquid properties
from CoolProp \citep{bell2014coolprop}.

\paragraph{Positioning.}
The ingredients exist: GRPO, verifiable rewards, solver-grounded reinforcement learning and language
models for process design. The combination for data-center energy systems, with second-law verifiers
as the reward, four task families with structural splits, an equation-level ledger benchmark and
accounting in solver calls, is what we contribute. \Cref{tab:positioning} summarizes how the cited
methods differ.

\begin{table}[tbp]
\centering\small\renewcommand{\arraystretch}{1.15}
\begin{tabular}{@{}p{3.4cm}p{1.9cm}p{2.4cm}p{2.4cm}p{1.7cm}p{1.7cm}@{}}
\toprule
Work & Language model & Trained with solver or verifier reward & Exergy or second-law checks & Data-center domain & Equation-level fault repair \\
\midrule
ThermoQA \citep{duzkar2026thermoqa} & evaluated & no & \todo{check} & no & no \\
Gas-path agents \citep{song2024gaspath} & prompted & no & \todo{check} & no & no \\
Text2Sim \citep{tian2026text2sim} & prompted, frozen & no & \todo{check} & no & no \\
Cycle co-design RL \citep{li2026thermocycle} & none (graph RL) & learned surrogate & \todo{check} & no & no \\
Solver-informed RL \citep{chen2025sirl} & trained & solver & \todo{check} & no & no \\
Nuclear TH modeling \citep{luo2026nuclear} & trained & solver & \todo{check} & no & no \\
\textbf{This work} (\sys) & to be trained & simulator and verifiers & yes (V3, V5) & yes & yes \\
\bottomrule
\end{tabular}
\caption{Positioning among cited work. Entries for prior work follow the descriptions recorded in
the project notes for each paper; ``no'' means the property is not part of the described method, and the
second-law column is not yet checked against the papers themselves. The last row describes the design;
the trained model is pending.}
\label{tab:positioning}
\end{table}
